\section{Related Work}
\label{sec:related}

\rev{
%Splitwise is specialized to the phased behavior of generative LLM inference.
%Several prior works on heterogeneous workloads and hardware in the field of scheduling.

\myparagraph{Heterogeneous scheduling and dataflow systems}
Prior work has studied heterogeneous scheduling for a variety of interactive services~\cite{zhang2019_MArk,patel2023hybrid,romero2019_INFaaS}.
These works exploit hardware heterogeneity to strike a balance between different objectives such as cost, energy, and performance.
However, they run the entire workload on the same machine.
Research on heterogeneous multiprocessor CPU scheduling attempts to match workload heterogeneity to hardware heterogeneity~\cite{haque2015few,van2012scheduling,haque2017fasttoslow,kumar2003single,yang2017powerchief,chen2009efficient}.
These works use profiling or online monitoring with metrics like request length or hardware performance counters to identify workload phases and allocate them appropriately on heterogeneous processors.
However, they do not consider the complexities with batching.
%
%Splitwise exploits the known two-phase structure of LLM inference requests and relies on profile-guided heuristics to support efficient batching for requests either in the same or in different phases. 
%\myparagraph{Distributed systems}
Distributed dataflow systems orchestrate large-scale computational graphs and aim to provide general-purpose programmability~\cite{dean2008mapreduce,isard2007dryad,shvachko2010hadoop,zaharia2012resilient}.
LLM inference under Splitwise can be viewed as a static computational graph with two stages, so it could be implemented using distributed frameworks that provide efficient GPU abstractions~\cite{moritz2018ray}.
Splitwise differs from these works since it uses a specialized two-phase design for generative LLM inference and leverages phase-aware resource management with efficient batching. 

\myparagraph{Model serving systems}
LLM inference serving is a rapidly developing field, with several recent works optimizing batching~\cite{yu2022orca,vllm,alpaserve,sarathi,aminabadi2022deepspeed}, scheduling~\cite{vllm,pope2023efficiently,sheng2023fairness,wu2023fast,turbomind,hong2023flashdecoding++}, and memory usage~\cite{flashattention,vllm,dettmers2022llmint8,sheng2023flexgen}.
Prior work has also proposed using CPUs and lower compute capability devices for LLM serving~\cite{bigdlllm,numenta}.
These approaches use the same machine for both prompt and token phase.
With Splitwise, they could improve throughput and latency by splitting phases.

%\myparagraph{Deep learning pipeline serving systems}
Prior work on video and ML serving focuses on scheduling model chains with data dependencies under latency constraints~\cite{hu2021scrooge,crankshaw2020inferline,romero2019_INFaaS,kannan2019_GrandSLAm,albahar2022schedtune}.
Such schedulers rely on model profiling to make efficient allocation decisions and manage requests across machines.
%However, unlike existing schedulers that orchestrate requests across models, Splitwise splits the phases of a single LLM inference request onto different machines via KV-cache transfer, and optimizes the complete system design after the split.
Recommendation system inference exhibits compute/memory heterogeneity both within and across models.
Prior work exploits such heterogeneity to selectively schedule requests between CPUs and accelerators~\cite{gupta2020deeprecsys,kwon2019tensordimm}, colocate models with complementary memory usage~\cite{choi2023hera}, and partition compute/memory on heterogeneous hardware resources~\cite{hwang2020centaur,jiang2021fleetrec}.
Similarly, Splitwise exploits the heterogeneity within LLM inference requests.
However, it uses different optimizations due to the differences in LLM workload characteristics and requirements.
%(\eg, recommendation systems requests are typically run to completion on a single machine due to smaller timescales and tighter latency constraints~\cite{gupta2020deeprecsys}).
}
